 \documentclass{article}

\title{Diploma}
\author{mishamilovidov}
\date{August 2025}

\usepackage{a4wide}
\usepackage[utf8]{inputenc}
\usepackage[T2A]{fontenc}
\usepackage[english,russian]{babel}
\usepackage{graphicx}
\usepackage{amsmath}
\usepackage{float}
\usepackage{hyperref}
\usepackage{amssymb}
\usepackage{bbm}
\usepackage{amsthm}
\usepackage{indentfirst}
\usepackage{float}
\usepackage{comment}
\usepackage[thinc]{esdiff}


\begin{document}

\thispagestyle{empty}

\begin{center}
\ \vspace{-3cm}


\includegraphics[width=0.5\textwidth]{msu.eps}\\
{\scshape Московский государственный университет имени М.~В.~Ломоносова}\\
Факультет вычислительной математики и кибернетики\\
Кафедра системного анализа

\vfill

{\LARGE Лабораторная работа по курсу ППП}

\vspace{1cm}

{\Huge\bfseries <<Задача о поиске места для заповедника>>}
\end{center}

\vspace{1cm}

\begin{flushright}
  \large
  \textit{Студент 415 группы}\\
  М.\,В.~Миловидов

  \vspace{5mm}

  \textit{Руководитель работы}\\
  Паршиков М. В.
\end{flushright}

\vfill

\begin{center}
Москва, 2026
\end{center}
\newpage

\section{Постановка задачи}

Имеется территория, на которой расположены несколько населённых пунктов. В пределах этой же территории планируется создать заповедник. Поскольку близость к населённым пунктам негативно влияет на качество воздуха и снижает ценность территории для охраны природы, заповедник желательно разместить как можно дальше от них. Вы (как представитель природоохранной организации) решаете выбрать оптимальное положение заповедника. При этом предполагается, что стоимость земли и её содержания почти постоянна и не зависит от выбранного места в пределах рассматриваемой территории. В качестве исходных данных необходимо использовать текстовый файл с координатами населённых пунктов (для каждого пункта в отдельной строке указаны широта и долгота его положения, через запятую). Широта и долгота задаются как вещественные числа, в градусах (без выделения минут и секунд); символ-разделитель — «точка». Также задаётся форма и границы общей территории, в пределах которой можно разместить заповедник. Основная задача состоит в разработке алгоритма поиска оптимального положения заповедника. При этом необходимо рассмотреть несколько разных постановок задачи, связанных с формой общей территории и приоритетами населённых пунктов: 
\begin{enumerate}
    \item Все населённые пункты равнозначны. Можно считать, что Земля является «плоской», так как расстояния между разными пунктами невелики. В качестве расстояния от заповедника до населённого пункта можно использовать расстояние по прямой. Общая территория, в пределах которой расположены населённые пункты и может быть создан заповедник, представляет собой прямоугольник, заданный своими границами. Требуется выбрать точку внутри этого прямоугольника для создания заповедника так, чтобы она находилась как можно дальше от всех населённых пунктов.
    \item Населённые пункты неравнозначны: для каждого указан приоритет (чем больше число, тем важнее, чтобы заповедник находился далеко от этого пункта, например из‑за большей численности населения или хозяйственной активности). Как и в первом пункте, Земля считается «плоской», расстояние измеряется по прямой. Общая территория может быть как прямоугольной, так и круглой. Требуется выбрать точку внутри этой территории для создания заповедника так, чтобы она находилась как можно дальше от населённых пунктов, особенно от тех, у которых высокий приоритет.
\end{enumerate}
Для каждого варианта задачи необходимо придумать постановку соответствующей математической задачи оптимизации и алгоритм численного решения этой задачи. Результатом выполнения работы являются программа и краткий отчёт. Программа должна выдавать оптимальные координаты положения заповедника.

Пусть \(n\) --- число населенных пунктов, \(X = [x_i], Y = [y_i] \in \mathbb{R}^n\) --- координаты населенных пунктов, \(W = [w_i] \in \mathbb{R}^n\) ---  веса городов (в случае первой постановки задачи будем считать \(w_i = 1,\ i = \overline{1, n}\)), \(\varepsilon\) ---  некая малая константа. 
Определим несколько функций, которые будем максимизировать в ходе работы:
\[f_1(x, y) = \sum_{i = 1}^n w_i\sqrt{(x - x_i)^2 + (y - y_i)^2},\]
\[f_2(x, y) = \sum_{i = 1}^n w_i((x - x_i)^2 + (y - y_i)^2),\]
\[f_3(x, y) = \min_{i = \overline{1, n}}\{\frac{\sqrt{(x - x_i)^2 + (y - y_i)^2}\}}{w_i + \varepsilon},\]
\[f_4(x, y) = \frac{1}{\sum_{i = 1}^n \frac{w_i}{\sqrt{(x - x_i)^2 + (y - y_i)^2}}}.\]
Они определяют разные подходы к определению оптимального положения заповедника.

\section{Решение}
Для поиска оптимальной точки будем использовать метод имитации отжига (simulated annealing) с некоторыми модификациями. Определим \(k\) --- коэффициент охлаждения \(T_i = kT_{i - 1}\) --- температура. Алгоритм стартует в случайной точке в области, далее действует следующий цикл: генерируется новая точка из нормального распределения (\(x^i \sim N(x^{i-1}, \text{stepscale}),\ y^i \sim N(y^{i-1}, \text{stepscale})\)), где \(\text{stepscale} = \text{step}_0 * (T_i / T_0)\). Если значение оптимизируемой функции в этой точке больше, то переход совершается гарантированно, иначе --- с вероятностью \(e^{\frac{\Delta f}{T + \varepsilon}}\), где \(\Delta f = f(x^i, y^i)\). Далее происходит локальная оптимизация --- проверяется 8 точек вокруг выбранной на расстоянии \(\text{localstep} = \frac{\text{stepscale}}{2}\). Локальный шаг уменьшается, пока не будет найдена более оптимальная точка или не будет достигнут минимальный локальный шаг. Далее алгоритм переходит к следующей итерации. Алгоритм завершается, когда будет достигнут минимальный локальный шаг.

\section{Примеры}
Везде будем использовать следующие параметры:

Для прямоугольника:
\[x_{max} = 100,\ x_{min} = -100,\ y_{max} = 100,\ y_{min} = -100.\]
Для круга:
\[x_c = 0,\ y_c = 0,\ r = 100.\]
\(x_{max},\ x_{min},\ y_{max},\ y_{min}\) определяются как углы описанного квадрата.
Число населенных пунктов \(N = 100\).

\[T_{min} - 10^{-8}, \varepsilon = 10^{-12}\]
\[T_0 = \frac{\sqrt{(x_{max} - x_{min})^2 + (y_{max} - y_{min})^2}}{4},\ \text{step}_0 = \max(x_{max} - x_{min}, y_{max} - y_{min})\]

\begin{figure}[H]
\includegraphics[width=1\textwidth]{data/ex2.eps}
\caption{Пример 1: прямоугольник, случайные города, функция \(f_1\)}
\end{figure}

\begin{figure}[H]
\includegraphics[width=1\textwidth]{data/ex3.eps}
\caption{Пример 2: прямоугольник, случайные города, функция \(f_2\)}
\end{figure}

\begin{figure}[H]
\includegraphics[width=1\textwidth]{data/ex4.eps}
\caption{Пример 3: прямоугольник, случайные города, функция \(f_3\)}
\end{figure}

\begin{figure}[H]
\includegraphics[width=1\textwidth]{data/ex5.eps}
\caption{Пример 4: прямоугольник, случайные города, функция \(f_4\)}
\end{figure}

\begin{figure}[H]
\includegraphics[width=1\textwidth]{data/ex1.eps}
\caption{Пример 5: прямоугольник, сетка городов с одним отсутствующим, функция \(f_3\)}
\end{figure}

\begin{figure}[H]
\includegraphics[width=1\textwidth]{data/ex6.eps}
\caption{Пример 6: круг, случайные города, функция \(f_3\)}
\end{figure}

Как видно, наиболее адекватный результат в данном случае показывает функция \(f_3\), то есть расстояние до взвешенного ближайшего населенного пункта. Максимизируя его, мы получаем точку, максимально удаленную от всех городов.

\end{document}
